\documentclass[a4paper]{article}
\usepackage[utf8]{inputenc}
\usepackage[T1]{fontenc}
\usepackage[french]{babel}
\usepackage{graphicx}
\usepackage{hyperref}
\usepackage{amsmath,amssymb}
\usepackage{geometry}
\geometry{margin=2.5cm}
\title{Projet d'électronique E3 EISC - Partie 1}
\author{Votre Nom}
\date{2025-2026}
\begin{document}
\maketitle


1. Le Contexte : Une PME en mutationL'entreprise : Jean Perret, PME familiale de 100 salariés (CA de 7,26 M€) située en Haute-Savoie.Le problème : Son activité historique (sous-traitance automobile) est peu rentable (1,5% de marge) en raison d'une forte concurrence.L'innovation : En répondant à une commande spécifique du client Sidel, l'entreprise invente un amortisseur hydraulique révolutionnaire (laminage de fluide avec membrane air/huile) plus performant et durable que l'existant.2. Opportunité et StratégiePropriété Industrielle : Un brevet national est déposé et obtenu fin $n-1$.Marché : Le PDG identifie un fort potentiel dans toute l'industrie (convoyeurs, presses, machines-outils).Le Projet : Transformer ce prototype en une activité stratégique propre, moins dépendante de la sous-traitance et plus rentable.3. L'Étude de Faisabilité (Avant-projet)Équipe : Le projet est confié à un stagiaire ingénieur (futur chef de projet) sous la responsabilité du PDG.Résultats :Technique : Maîtrise des technologies de fabrication en série confirmée.Économique : Excellente rentabilité prévue, malgré une menace de concurrents développant d'autres principes.Planification initiale : Durée évaluée à 8 mois pour un coût de 126 K€.4. Organisation et FinancementAide de l'ANVAR : L'agence finance 75% du coût du projet via une avance sans intérêts, remboursable sur 4 ans dès le début de l'exploitation.Ressources : Le projet mobilise diverses expertises internes (Bureau d'études, Méthodes, Achats, Production) à temps partiel, sauf le chef de projet qui est à 100%.Calendrier : Lancement officiel le 7 juillet $n+1$, avec des contraintes de congés et de jours fériés (base de 20 jours travaillés par mois).5. Aléas et Pilotage (La réalité du terrain)Le conflit de ressources : Une surcharge imprévue des ateliers (commandes de sous-traitance) menace de retarder la fabrication des prototypes de 8 semaines.Le dilemme du manager :Option A : Accepter le retard (risque de rater la fenêtre de marché face aux concurrents).Option B : Sous-traiter la fabrication des prototypes pour 23 K€ supplémentaires afin de tenir les délais.Points clés à retenir pour ton cours :Transition stratégique : Passer de la sous-traitance (faible valeur ajoutée) à une activité de fabricant (produit propre, breveté).Gestion des contraintes : Équilibrer le planning entre les ressources humaines limitées et les impératifs de production.Arbitrage financier : Le coût de la sous-traitance (23 K€) vs le coût d'opportunité d'un retard de deux mois.



% --- Fin du document ---
\end{document}
